\begin{table}[htbp]
\centering
\footnotesize
\caption{Resultados do Modelo BMA Principal — Desfecho: Letalidade}
\label{tab:bma_letalidade}
\begin{tabular}{lccccr}
\toprule
\textbf{Covariável} & \textbf{PIP} & \textbf{Média Post.} & \textbf{DP Post.} & \textbf{P(Sinal +)} & \textbf{Direção} \\
\midrule
Município integrante de região metropolitana & \textbf{0,7809} & \textbf{0,4517} & 0,3372 & 1,0000 & Positiva \\
Proporção de mulheres eleitas & \textbf{0,7441} & \textbf{-0,0937} & 0,0788 & 0,0000 & Negativa \\
Cobertura média de imunização em menores de um ano & \textbf{0,6355} & \textbf{-0,0598} & 0,0678 & 0,0000 & Negativa \\
Alinhamento partidário do prefeito com o governo federal & \textbf{0,6161} & \textbf{0,1116} & 0,1277 & 1,0000 & Positiva \\
IDEB dos anos finais do ensino fundamental & \textbf{0,5656} & \textbf{-0,0522} & 0,0734 & 0,0000 & Negativa \\
IVS Capital Humano & \textbf{0,5644} & \textbf{0,0717} & 0,0926 & 1,0000 & Positiva \\
Leitos vinculados ao SUS & \textbf{0,5635} & \textbf{0,0561} & 0,0764 & 1,0000 & Positiva \\
Proporção da população em extrema pobreza & \textbf{0,5466} & \textbf{-0,0848} & 0,1174 & 0,0012 & Negativa \\
Internações por doenças do aparelho respiratório & \textbf{0,5365} & \textbf{-0,0529} & 0,0741 & 0,0001 & Negativa \\
Unidades de saúde vinculadas ao SUS & \textbf{0,5336} & \textbf{-0,0431} & 0,0670 & 0,0000 & Negativa \\
População beneficiária do Programa Bolsa Família & \textbf{0,5285} & \textbf{0,0564} & 0,0803 & 0,9989 & Positiva \\
Área territorial & 0,4824 & 0,0306 & 0,0552 & 0,9996 & Positiva \\
IDEB do 3º ano do ensino médio & 0,4813 & 0,0306 & 0,0539 & 0,9998 & Positiva \\
Internações por doenças do aparelho circulatório & 0,4749 & 0,0343 & 0,0654 & 0,9972 & Positiva \\
Proporção da população em empregos formais & 0,4675 & -0,0355 & 0,0734 & 0,0020 & Negativa \\
Taxa de homicídios & 0,4560 & -0,0234 & 0,0499 & 0,0112 & Negativa \\
Índice de Desenvolvimento Municipal & 0,4560 & -0,0396 & 0,0856 & 0,0210 & Negativa \\
População estimada & 0,4538 & 0,0208 & 0,1401 & 0,8219 & Positiva \\
Índice de Desenvolvimento Humano Municipal & 0,4527 & 0,0382 & 0,0975 & 0,9043 & Positiva \\
IVS Infraestrutura Urbana & 0,4496 & -0,0221 & 0,0524 & 0,0139 & Negativa \\
Densidade demográfica & 0,4461 & 0,0252 & 0,1400 & 0,8782 & Positiva \\
IDEB dos anos iniciais do ensino fundamental & 0,4309 & 0,0010 & 0,0572 & 0,6693 & Positiva \\
Consumo de energia elétrica por 100 mil habitantes & 0,4306 & -0,0194 & 0,0572 & 0,0136 & Negativa \\
Profissionais de saúde vinculados ao SUS & 0,4260 & -0,0120 & 0,0585 & 0,1556 & Negativa \\
Produto Interno Bruto per capita & 0,4231 & -0,0084 & 0,0600 & 0,3253 & Negativa \\
Proporção da população rural & 0,4169 & -0,0053 & 0,0523 & 0,4223 & Negativa \\
Proporção da população com 60 anos ou mais & 0,4159 & 0,0091 & 0,0483 & 0,8976 & Positiva \\
IVS Renda e Trabalho & 0,4134 & 0,0101 & 0,0646 & 0,8429 & Positiva \\
Proporção de votos válidos do prefeito eleito & 0,4124 & -0,0051 & 0,0378 & 0,2579 & Negativa \\
Despesa municipal com saúde e saneamento & 0,4105 & 0,0124 & 0,0556 & 0,8689 & Positiva \\
Município integrante do semiárido & 0,4085 & 0,0135 & 0,1014 & 0,8696 & Positiva \\
Despesa municipal com educação e cultura & 0,4061 & 0,0037 & 0,0454 & 0,7425 & Positiva \\
Internações por diabetes mellitus & 0,4054 & 0,0023 & 0,0476 & 0,6600 & Positiva \\
Internações por asma & 0,4053 & -0,0050 & 0,0436 & 0,2039 & Negativa \\
Cobertura urbana de abastecimento de água & 0,4012 & -0,0031 & 0,0357 & 0,1738 & Negativa \\
\bottomrule
\end{tabular}
\begin{tablenotes}\footnotesize
\item \textit{Nota:} Variáveis em negrito possuem $\text{PIP} \ge 0{,}50$. Estimativas obtidas via amostrador MCMC (2 cadeias de 1.000.000 de iterações).
\end{tablenotes}
\end{table}